\documentclass[a4paper]{article}
\usepackage[utf8]{inputenc}
\usepackage[T1]{fontenc}
\usepackage{amsmath}
\usepackage{amssymb,amsthm}
\usepackage{graphicx}
\usepackage{mathtools}
\usepackage[framemethod=tikz]{mdframed}
\usepackage{enumitem}
\usepackage{csquotes}
\usepackage{mathrsfs}
\usepackage{setspace}
\usepackage{xfrac}
\usepackage{scalerel}[2014/03/10]
\usepackage{stackengine}




\title{Master thesis notes}
\author{Elia Perli}


\renewcommand\thesection{\Alph{section}} %← mette le sezioni con le lettere anziché i numeri

\newtheorem{theo}{Theorem}
\surroundwithmdframed[outerlinewidth=0.4pt,
innerlinewidth=0.4pt,
middlelinewidth=1pt,
middlelinecolor=white,
bottomline=false,topline=false,rightline=false]{theo}

\renewcommand{\phi}{\varphi}
\renewcommand{\epsilon}{\varepsilon}
\renewcommand{\Omega}{\varOmega}
\renewcommand{\Gamma}{\varGamma}
\renewcommand{\iff}{\Leftrightarrow}
\renewcommand{\rho}{\varrho}
\renewcommand{\div}{\operatorname{div}}



\newcommand{\R}{\mathbb{R}}
\newcommand{\C}{\mathbb{C}}
\newcommand{\Z}{\mathbb{Z}}
\newcommand{\N}{\mathbb{N}}
\newcommand{\D}{\mathcal{D}}
\newcommand{\M}{\mathbb{M}^{2,1}}

\newcommand{\htl}{\reallywidetilde{h^{-1}}}
\newcommand{\Leb}{\mathscr{L}}
\newcommand{\restr}[1]{_{|_{#1}}}
\newcommand{\sgn}{\operatorname{sgn}}
\newcommand{\then}{\Rightarrow}
\newcommand{\supp}{\operatorname{supp}}
\newcommand{\Dom}{\operatorname{Dom}}


\newcommand{\comment}[1]{\text{\parbox{5cm}{\setstretch{0.1}#1}}}
\newcommand{\astr}[1]{\text{$(*#1)$}}

\newtheorem{manualtheoreminner}{\textsc{teorema}}
\newenvironment{manualtheorem}[1]{%
	\renewcommand\themanualtheoreminner{#1}%
	\manualtheoreminner
}{\endmanualtheoreminner}

\newtheorem{thm}{\textsc{Theorem}}[section]
\newtheorem{cor}{\textsc{Corollary}}[section] %,\Alph] %[section/chapter, ...] \newtheorem{prop}{\textsc{Proposition}}[section]
\newtheorem*{prop*}{\textsc{Proposition}}
\newtheorem{prop}{\textsc{Proposition}}
\newtheorem{oss}{\textsc{remark}}[section]
\theoremstyle{definition}
\newtheorem{defn}{\textsc{Definition}}[section]
\newtheorem{ntt}{\textsc{Notation}}[section]
\newtheorem{es}{\textsc{example}}[section]



% figure support
\usepackage{import}
\usepackage{xifthen}
\pdfminorversion=7
\usepackage{pdfpages}
\usepackage{transparent}
\newcommand{\incfig}[1]{%
	\def\svgwidth{\columnwidth}
	\import{./figures/}{#1.pdf_tex}
}

\pdfsuppresswarningpagegroup=1

\begin{document}

\tableofcontents



\section{Introduction}
Goal of the work, maybe a simple example, some mention to preliminary information and a soft mention to what's the problem and the philosophy.
\section{Minkowski space and space-like surfaces}
The environment we will work within is the Minkowki 3-space, defined as $\R^3$ with the pseudometric
\[ds^2 = ( dx^1)^{2} + (dx^2)^{2}-(dx^3)^2.\]
In order to distinguish this space from $\R^3$ we will use the notation $\M$. We will also denote the coordinates as $(x^1,x^2,x^3)$ or equivalently $(x,y,t)$.

\paragraph{Space-like immersions.}\mbox{}\\
The object we will focus on are the so-called space-like surfaces of $\M$. This term refers to the immersed surfaces in $\M$ whose induced metric is positive-definite hence Riemannian.\\

A nice property of this kind of surface is that they are locally the graph of a function of the form $x^3 = f(x^1,x^2)$.\\
This fact can be seen with a simple reasoning:\\
First, we assumed that the pseudometric restricted to $S$,$ \left< \cdot , \cdot \right>|_S \geq 0 $, is definite positive, which means that the tangent vectors of $S$ at any point $p$ belong to the set of space-like vectors:
\begin{equation*}
	T_p \subset \mathscr{C} := \left\{ v\in \M : \left<v,v \right> \ge 0\right\}=\{ t^2 \leq  x^2+y^2\} . 
\end{equation*}
This is just the outside-region of a double cone in $\M$, whose symmetry axis is $ \left\{ x^1=0, x^2 = 0 \right\} $. A plane is entirely made of space-like vectors iff it's contained between the two branches of the cone.\\
Second, having assumed that $S$ is an immersed surface, hence an immersed submanifold, it is locally the solution of an equation of the form  $F(x,y,z)=0$. It is well known that from this local expression we can express the normal vector to the surface at a point $p$: $\vec n_p = \vec \nabla F(p)$.\\
Furthermore, we observe that, being $\vec n_p$ orthogonal to $T_pS$, it cannot be horizontal,  $(\vec n_p)^3 = \partial_z F(p) \neq 0$, which means that we can apply the implicit function theorem and express the surface as the graph of $z=f(x,y)$.\\
Furthermore, where the local representation holds, we can parametrize the surface as $(x,y,f(x,y))$, from which two generators of the tangent space are $V_x = (1,0,\partial_x f(x,y))$ and $V_y = (1,0,\partial_y f(x,y))$, and the first fundamental form is represented by the matrix
\begin{align*}
	g &= 
\begin{pmatrix}
\left<V_x,V_x \right> & \left<V_y,V_x \right>\\
\left<V_x,V_y \right> & \left<V_y,V_y \right>
\end{pmatrix}
	  &=
\begin{pmatrix}
1-f_x^2 & -f_xf_y\\
-f_xf_y & 1-f_y^2
\end{pmatrix}.
\end{align*}
The condition on this being a metric translates in having a strictly positive determinant  (as the parametrization is regular).
\begin{align*}
	0< \det g  = 1-(f_x^2+f_y^2) = 1-|\vec \nabla f|^2,
\end{align*}
which means $|\vec \nabla f|^2 <1$.
It's not difficult to convince ourselves, as the local representation as a graph is always possible, that the conditions of being able to explicit the time variable as dependent and that $|\vec \nabla f|^2 <1$ are a characterization of space-like immersed surfaces.\textcolor{red}{mettere dopo}\\
\paragraph{Entire surfaces.}\mbox{}\\
Given a surface $S$ immersed in  $\M$, we can define a projection $\Pi:S\to \R^2$ by forgetting the time coordinate, $\Pi(x^1,x^2,x^3):=(x^1,x^2)$ for any $(x^1,x^2,x^3)\in S$. We call the surface $S$ is entire when the map $\Pi$ is surjective. If the surface is globally a graph of a function $\phi$, this condition is exactly stating that the function $\phi$ is an entire map from $\R^2$ to $\R$.\\

The situation for space-like surfaces is actually a lot nicer than what we presented up to now, as the representation as a graph is not only local but actually global, as follows from the proposition below. \textcolor{red}{dovrei mostrare che davvero la rappresentazione è globale?}
\begin{prop}
	Let $\Omega$ be a domain in $\C$ (that is, an open connected set) and let $S:\Omega \to \M$ be a space-like immersion. Suppose that $S$ is complete with respect to the induced Riemannian metric. Then $S$ is an entire graph and $\Omega$ is simply connected.
\end{prop}
\paragraph{Immersion data for $S$}\mbox{}\\
First, second fundamental forms and curvatures
\paragraph{The Gauss map in $\M$}\mbox{}\\
The target of the gauss map is $\mathbb{H}^2$ and why, and it's isometric to the Poincaré Disk.

\begin{itemize}[noitemsep]
	\item $\mathbf{M}^{2,1}$
	\item Space-like immersion, induced Riemannian metric
	\item First and second fundamental forms
	\item Mean curvature and Gaussian curvature expr
	\item Space-like surfaces are graphs of functions
	\item The Gauß map has  $\mathbb{H}^2$ as a target
\end{itemize}

\section{Harmonic maps}
\section{Classification of CMC}
\subsection{The idea of the proof}
\subsection{The analytical work}
\end{document}




\begin{thm}
	Let $ h$ be a continuous function defined on a domain $\Omega$ in $\C$. Let $\tilde w$ be a solution for [where do we use the inequality with zero?]
	\begin{align*}
		\Delta\tilde{w} &= e^{2\tilde{w}} - h^{2} e^{-2\tilde{w}} \ge 0
	\end{align*}
	in  $\Omega$ such that $e^{2\tilde w}$ is a complete metric on $\Omega$. Then any other solution $v$ of the equation satisfies $v\leq \tilde w$.
\end{thm}
\begin{proof}
	
\end{proof}

\begin{cor}{}{}
	Let $h$ be a continuous function defined on a domain $\Omega\subset \C$ and let $\tilde w_1, \tilde w_2$ two solutions of 
	\begin{equation}
		\Delta \tilde w  = e^{2\tilde w}-h^2e^{-2\tilde w} \geq 0
	\end{equation}
	in $\Omega$ such that $e^{2\tilde w_i}$ are complete metrics on $\Omega$. Then  $\tilde w_1 \equiv \tilde w_2$.
\end{cor}
\begin{cor}
	There is at most one solution of the equation ??? such that the associated metric is complete.
\end{cor}

